% ---- pg-024 （书页 12） ----

\[
\begin{aligned}
x_1 &= \text{密度(克 / 立方厘米)}\\
x_2 &= \text{强度(磅): 沿机器方向}\\
x_3 &= \text{强度(磅): 沿交叉方向}
\end{aligned}
\]

一种新颖的图形用来表示这些数据,见\hyperref[fig:1.5]{图 1.5}.散布图作为协方差阵列的非对角线元素,而盒形图作为对角线元素.用这一方法后者有不同的标度,故我们只用总的形状对称性给出信息,并为每个单独的特性提供可能的离群值.可对图形和异常的观测值来观察散布图.在\hyperref[fig:1.5]{图 1.5} 中,有一个异常的观测值: 样本 25 的密度.某些散布图表明,其观测值是相互分离的两簇.阵列的这些散布图将在下一节中新的软件图解方法的讨论中还会涉及.

\begin{mvfigure}
\mvinc{pg-024}{0.13938,0.28892,0.84977,0.64921}
\phantomsection\label{fig:1.5}\mvcap{图 1.5\quad 来自表 1.2 的纸张性能数据的散布图和盒形图}
\end{mvfigure}

在一般的多种响应的情况下, $p$ 个变量同时进行 $n$ 项记录.应该为几对重要的变量作散布图,并且如果工作量没有大到得不偿失,应该对每一对变量作图.

由于我们被限定在一个三维世界,我们经常无法描述整个的数据集合.然而,两种进一步的数据几何表示为观察多元统计方法提供了一个重要的概念框架.假如在三维空间捕捉到数据的本质是可能的,就可以用实际图形表示.

\mvbold{$p$ 维空间中的 $n$ 个点($p$ 维散布图)}\quad 考虑将散布图自然地扩展到 $p$ 维空间,其中 $p$ 个观测结果

\[
(x_{j1}, x_{j2}, \cdots, x_{jp})
\]

表示一个点在 $p$ 维空间中的坐标在第 $j$ 项的值.每个坐标轴与变量相对应,则第 $j$ 个点沿着第一个坐标轴是 $x_{j1}$ 个单位,沿第二个坐标轴是 $x_{j2}$ 个单位,$\cdots$,沿第 $p$ 个坐标轴是 $x_{jp}$ 个单位的点.由此得到的有 $n$ 个点的结果的图形不仅显示了变异性的总的模型,而且会显示出 $n$ 项中
